\documentclass{ipgplicence}

\begin{document}


\status{Received}
%\status{Corrected}
%\status{Accepted}

\selectlanguage{french}
%\selectlanguage{english}

% Éviter les "\\" pour séparer les paragraphes (insérer plutôt une ligne blanche)
% Éviter de forcer le positionnement des figures avec l'option de "[h!]"

%\setkeys{Gin}{draft} % images non chargées -> espaces réservés et compilation plus rapide
\setkeys{Gin}{draft=false} % images chargées

\thispagestyle{empty}

% Insérer ici les différents logos que vous voulez faire appaître
\includegraphics[width=0.28\linewidth]{figures/LogoIPGP_UPC}\hfill%
\includegraphics[width=0.25\linewidth]{figures/p7_logo}\hfill%
\includegraphics[width=0.28\linewidth]{figures/LogoIPGP_UPC}%

\vspace{1cm}

\begin{center}
{\LARGE\bfseries
Titre du Rapport
}

\vspace{2mm}

{\Large\bfseries
Nom Prénom$^1$,  Nom Prénom$^2$, Nom Prénom$^3$...
}

\vspace{4mm}

$^1$ Institution1, $^2$ Institution2, $^3$ Institution3

{\Large
Université d'Enseignement (si nécessaire)

\vspace{5mm}

Stage de \textit{LEVEL} effectué à \textit{WHERE} sous la direction de \textit{TUTOR} \\[2mm]
du \textit{BEGIN\_DATE} au \textit{END\_DATE}
}

\end{center}

\vspace{2mm}

\centerline{\today}
%\centerline{2 mai 2020}

\vspace{1cm}

\addcontentsline{toc}{section}{Résumé (en français)}
\subsection*{\centerline{Résumé (en français)}}

Un court résumé portant sur la problématique du projet et les résultats principaux obtenus, éventuellement en 2 langues.

\vspace{5mm}

\subsection*{\centerline{Résumé dans une autre langue, si nécessaire}}
\addcontentsline{toc}{section}{Résumé dans une autre langue, si nécessaire}

\includegraphics[width=\linewidth]{figures/japonais}


\newpage

{\parskip 0pt
\tableofcontents
}

\newpage

\section{Introduction}


À la lecture de l’introduction, le lecteur doit avoir compris  le cadre du sujet et le champ de vos investigations. 

L’introduction doit comporter :
\begin{itemize}
\item le sujet et les questions soulevées par celui-ci,
\item l’intérêt du sujet et de l’étude,
\item la délimitation du sujet : votre champ d’investigation, votre contribution et ce qui ne sera pas abordé,
\item la présentation succincte du plan (à ne pas oublier).
\end{itemize}

\section{Corps du rapport}

\begin{itemize}

\item Le corps du rapport doit être structuré (sections voire chapitres ou parties suivant la taille du document) et  contenir: 
\begin{enumerate}
\item une description du projet confié et des questions posées,
\item une présentation des méthodes utilisées,
\item une présentation et une analyse des résultats obtenus.
\end{enumerate}

\item Le nombre de sections peut varier suivant le type de sujet abordé dans le document et la forme du projet (succession de questions ou non).

\item Votre document est destiné à pouvoir être lu par différentes personnes. Il reflète votre travail. Il faut donc attacher une importance non seulement au fond mais aussi à la forme.

\item Le document doit être synthétique, clair, lisible et bien présenté (qualité des illustrations, de la mise en page, typographie, charte graphique, etc.). 

\item Il vous est recommandé d’utiliser modérément le ‘je’. 
\end{itemize}


	\section{Tables, figures et références croisées}

La table~\ref{tab:ici} représente un exemple de tableau, la figure~\ref{fig:la} d'une image. Toutes vos figures et vos tableaux doivent être accompagnés d'une légende et référencés dans le texte. Avec pdflatex, les figures insérées peuvent être au format pdf, png ou jpg.

\begin{table}[htb]
\caption{Une table a toujours une légende.}
\label{tab:ici} % sert à référencer la table dans le texte, à mettre juste après la caption
\begin{center}
\begin{tabular}{rrr} \toprule
colonne 1 & colonne 2 & colonne 3 \\ \midrule
1 & 2 & 3 \\
24 & 125 & NaN \\ \bottomrule
\end{tabular}
\end{center}
\end{table}

\begin{figure}[hbt]
\centering
\includegraphics[width=0.8\textwidth]{figures/bmp6-3.pdf}
\caption{Carte de la population mondiale répartie par pays. Source des données \cite{NE} \textit{Une figure aussi doit avoir une légende!}}
\label{fig:la}
\end{figure}

\section{Équations}
Les équations importantes doivent être numérotées afin que l'on puisse y faire référence dans le texte. En physique on a tendance à toujours numéroter en math c'est moins le cas. Les maths en ligne s'écrivent ainsi : $x=\cos(y)$.

La même équation sur une ligne seule mais non numérotée

$$
x = \cos(y) .
$$

Enfin votre équation préférée, numérotée car importante !
\begin{equation}
\Delta\phi = 0 .
\label{eq:laplace}
\end{equation}

On peut aussi utiliser les références croisées pour citer  (\ref{eq:laplace}) dans laquelle vous aurez reconnu l'équation de Laplace n'est-ce pas ?

\section{Codes informatiques}
Vous avez deux façons de saisir des codes informatiques. La plus simple est d'utiliser l'environnement \texttt{verbatim} comme ci-dessous qui permet de formater les lignes par soi-même

\begin{verbatim}

import numpy as np
import matplotlib.pyplot as plt

def plot_cos( min = 0, max = np.pi ):
    x = np.linspace( min, max, 100 )
    y = np.cos(x)

    fig = plt.figure()
    ax = fig.add_subplot( 111 )
    ax.plot(x,y,'r-')

    plt.show()
\end{verbatim}


Cependant dès que votre code devient plus complexe il est plus intéressant d'utiliser des paquets spécifiques comme \texttt{listings} qui donnent des résultats nettement plus lisibles.


%\begin{minted}[linenos,breaklines]{python}
\begin{lstlisting}[language=python]
import numpy as np
import matplotlib.pyplot as plt

def plot_cos( min = 0, max = np.pi ):
    x = np.linspace( min, max, 100 )
    y = np.cos(x)

    fig = plt.figure()
    ax = fig.add_subplot( 111 )
    ax.plot(x,y,'r-')

    plt.show()
\end{lstlisting} 
%\end{minted}

Et ce d'autant que \texttt{listings} permet d'utiliser plusieurs langages comme Python mais aussi le SQL :
%\begin{minted}[linenos,breaklines]{SQL}
\begin{lstlisting}[language=sql]
select anneeUFR, count(*) from IdentiteTbl group by anneeUFR
\end{lstlisting}
%\end{minted}

\section{Compétences et connaissances acquises}

\section{Conclusion}

Celle-ci doit reprendre l’objet du projet et les principaux résultats obtenus et éventuellement ce qui reste à traiter (dans le cas où la durée du stage est trop courte par rapport aux objectifs fixés). Elle porte également des suggestions sur les orientations à court terme par rapport au sujet traité. 

\section{Remerciements}

\section{Bibliographie}

La bibliographie est importante  et ne doit pas être négligée. Peuvent y figurer ouvrages, traités, manuels, certains périodiques ou revues spécialisées. Il ne faut citer que les documents sur lesquels vous vous êtes appuyés pour rédiger vos travaux. La description de la bibliographie doit permettre au lecteur de retrouver les documents référencés, et de les identifier. Les références doivent donc être précises et complètes. On citera les référence au moyen de \cite{Ref1} ou encore de \cite{Ref2,Ref3}.


Ces références seront regroupées en fin de texte et classées par ordre alphabétique d'auteurs et pour un auteur, par date croissante. Des exemples sont donnés ci-dessous.

Les bibliographies peuvent être générées soit manuellement, comme ici, soit avec le logiciel \textsc{Bib}\TeX, ou sa version plus récente \textsc{Bib}\LaTeX, ce qui est vivement conseillé. En effet, \textsc{Bib}*\TeX{} permet de maintenir une base de données pérenne et réutilisable des références bibliographiques, et permet de choisir l’apparence de la sortie, sans avoir à modifier sa base de données (fichier .bib). Pour plus de renseignements, voir, par exemple, la section “Bibliographies”  de la documentation
\url{http://www.ipgp.fr/~moguilny/LaTeX/ltxdoc.pdf}.



\begin{thebibliography}{1}
\addcontentsline{toc}{section}{Références}



%Articles de périodiques
\bibitem{Ref1} NOM, Prénom. Titre de l’article. \textit{Titre du périodique}, année de publication, \textbf{volume}, numéro, pagination. 

%Ouvrages
\bibitem{Ref2} NOM, Prénom. \textit{Titre de l’ouvrage}. Tomaison. Édition. Lieu d’édition : éditeur commercial, année de publication, nombre de pages (Titre de la collection, n° de la collection).

%Thèse, mémoire de master, rapport de stage
\bibitem{Ref3} NOM, Prénom. \textit{Titre de la thèse}. Discipline. Lieu de soutenance : Établissement de Soutenance, année de soutenance, nombre de pages

%Sites internet
\bibitem{Ref4} AUTEUR ou ORGANISME (responsable du site). Titre de la page d’accueil [en ligne]. (date de création du site, date de mise à jour)Disponible sur : <URL> (date de consultation).

\bibitem{NE} Natural Earth, Admin - 0 Countries. \href{https://www.naturalearthdata.com/downloads/10m-cultural-vectors/10m-admin-0-countries/}{https://www.naturalearthdata.com/downloads/10m-cultural-vectors/10m-admin-0-countries/} (09/2018)

\end{thebibliography}

\appendix

\section{Annexes}

Vous pouvez si nécessaire (pour ne pas alourdir votre document), renvoyer à des annexes (par exemple cf. annexe X) numérotées (de manière alphabétique de préférence) et paginées en fin de document. Les annexes sont particulièrement utiles si vous voulez détailler un code informatique, un protocole expérimental dont la description serait trop longue dans le corps du rapport, ou insérer votre convention de stage.

\newpage

\subsection{Exemple d'inclusion de Convention de Stage}

\includegraphics[width=0.9\linewidth]{figures/ConvStage}

\end{document}
